\documentclass[10pt,a4paper]{amsart}
\usepackage{amsmath,amssymb,amsthm,graphicx}
% Retained minimal preamble for the Stacks Project chapter more-algebra.tex (More on Algebra).
% The chapter file pinned-original-source.tex carries only \input{preamble} on its first line, so no
% preamble text lives inside the chapter itself.  The Stacks Project uses one single shared preamble for
% every chapter; the snapshot stages it identically as sources/stacks-chow/preamble.tex and
% sources/stacks-etale-cohomology/preamble.tex.  The block below is copied mechanically, line for line,
% from sources/stacks-chow/preamble.tex lines 14-18 (the Xy-pic setup) and 154-217 (the native theorem
% declarations together with the mathematical macros).  Nothing was added, removed or redefined; the
% wrapper class and the amsmath/amssymb/amsthm/graphicx base packages are supplied by the assembler.
% Chapter: more-algebra.tex -- More on Algebra.
% For commutative diagrams we use Xy-pic
\usepackage[all]{xy}

% We use 2cell for 2-commutative diagrams.
\xyoption{2cell}

% Theorem environments.
%
\theoremstyle{plain}
\newtheorem{theorem}[subsection]{Theorem}
\newtheorem{proposition}[subsection]{Proposition}
\newtheorem{lemma}[subsection]{Lemma}

\theoremstyle{definition}
\newtheorem{definition}[subsection]{Definition}
\newtheorem{example}[subsection]{Example}
\newtheorem{exercise}[subsection]{Exercise}
\newtheorem{situation}[subsection]{Situation}

\theoremstyle{remark}
\newtheorem{remark}[subsection]{Remark}
\newtheorem{remarks}[subsection]{Remarks}

\numberwithin{equation}{subsection}

% Macros
%
\def\lim{\mathop{\mathrm{lim}}\nolimits}
\def\colim{\mathop{\mathrm{colim}}\nolimits}
\def\Spec{\mathop{\mathrm{Spec}}}
\def\Hom{\mathop{\mathrm{Hom}}\nolimits}
\def\Ext{\mathop{\mathrm{Ext}}\nolimits}
\def\SheafHom{\mathop{\mathcal{H}\!\mathit{om}}\nolimits}
\def\SheafExt{\mathop{\mathcal{E}\!\mathit{xt}}\nolimits}
\def\Sch{\mathit{Sch}}
\def\Mor{\mathop{\mathrm{Mor}}\nolimits}
\def\Ob{\mathop{\mathrm{Ob}}\nolimits}
\def\Sh{\mathop{\mathit{Sh}}\nolimits}
\def\NL{\mathop{N\!L}\nolimits}
\def\CH{\mathop{\mathrm{CH}}\nolimits}
\def\proetale{{pro\text{-}\acute{e}tale}}
\def\etale{{\acute{e}tale}}
\def\QCoh{\mathit{QCoh}}
\def\Ker{\mathop{\mathrm{Ker}}}
\def\Im{\mathop{\mathrm{Im}}}
\def\Coker{\mathop{\mathrm{Coker}}}
\def\Coim{\mathop{\mathrm{Coim}}}

% Boxtimes
%
\DeclareMathSymbol{\boxtimes}{\mathbin}{AMSa}{"02}

%
% Macros for moduli stacks/spaces
%
\def\QCohstack{\mathcal{QC}\!\mathit{oh}}
\def\Cohstack{\mathcal{C}\!\mathit{oh}}
\def\Spacesstack{\mathcal{S}\!\mathit{paces}}
\def\Quotfunctor{\mathrm{Quot}}
\def\Hilbfunctor{\mathrm{Hilb}}
\def\Curvesstack{\mathcal{C}\!\mathit{urves}}
\def\Polarizedstack{\mathcal{P}\!\mathit{olarized}}
\def\Complexesstack{\mathcal{C}\!\mathit{omplexes}}
% \Pic is the operator that assigns to X its picard group, usage \Pic(X)
% \Picardstack_{X/B} denotes the Picard stack of X over B
% \Picardfunctor_{X/B} denotes the Picard functor of X over B
\def\Pic{\mathop{\mathrm{Pic}}\nolimits}
\def\Picardstack{\mathcal{P}\!\mathit{ic}}
\def\Picardfunctor{\mathrm{Pic}}
\def\Deformationcategory{\mathcal{D}\!\mathit{ef}}

% --- appended 2026-10-12 (Chief): the project own definitions of the reference/slogan/history
% display environments, copied mechanically from the shared Stacks preamble lines 8-12
% (\usepackage{verbatim} plus the three \newenvironment lines, which the chapter preamble keeps
% verbatim from the project preamble). Nothing added, removed or redefined: these environments
% carry no mathematical content and the project itself typesets them as comments. Without these
% three lines a counted statement whose own text uses them fails with Environment slogan undefined.
\usepackage{verbatim}
\newenvironment{reference}{\comment}{\endcomment}
\newenvironment{slogan}{\comment}{\endcomment}
\newenvironment{history}{\comment}{\endcomment}
\usepackage{cleveref}
\crefname{originallemma}{original lemma}{original lemmas}
\Crefname{originallemma}{Original lemma}{Original lemmas}
\makeatletter
\newlabel{more-algebra-lemma-evaluate-and-more}{{0A60}{source}{}{sourcepointer-1}{}}
\newlabel{more-algebra-lemma-evaluate-and-more@cref}{{[originallemma][0][]0A60}{[1][0][]source}{}{}{}}
\makeatother
\crefname{originallemma}{original lemma}{original lemmas}
\Crefname{originallemma}{Original lemma}{Original lemmas}
\makeatletter
\newlabel{remark-internal-hom}{{0GWQ}{source}{}{sourcepointer-2}{}}
\newlabel{remark-internal-hom@cref}{{[originallemma][0][]0GWQ}{[1][0][]source}{}{}{}}
\makeatother
\crefname{originallemma}{original lemma}{original lemmas}
\Crefname{originallemma}{Original lemma}{Original lemmas}
\makeatletter
\newlabel{categories-remark-internal-hom-monoidal}{{0GWJ}{source}{}{sourcepointer-3}{}}
\newlabel{categories-remark-internal-hom-monoidal@cref}{{[originallemma][0][]0GWJ}{[1][0][]source}{}{}{}}
\makeatother
\title{A canonical morphism from a total tensor-internal-hom complex to an internal hom complex}
\author{The Stacks Project Authors}
\date{stacks-project:more-algebra.tex:local-snapshot-20261012; GNU FDL 1.2 or later}
\begin{document}
\maketitle
\setcounter{subsection}{1}
\noindent\textit{Editorial note.} All delivered mathematical text is the author's own, copied from the pinned Stacks Project chapter file more-algebra.tex; no statement or proof step was rewritten, reordered or removed.\par\medskip
\noindent\textit{Editorial note.} The chapter file pinned-original-source.tex carries only \textbackslash{}input\{preamble\} on its first line, so no preamble text lives inside the chapter. The retained file sources/stacks-more-algebra/retained-preamble.tex is a mechanical, line-for-line copy of the project's single shared preamble as staged in the snapshot (Xy-pic setup plus theorem declarations and mathematical macros). Nothing was added, removed or redefined, and the excerpt serves only as preamble, never as a span coordinate.\par\medskip
\noindent\textit{Editorial note.} Both counted spans start at the \textbackslash{}begin line of their environment, so the statement and the author's proof are delivered whole; surrounding prose of the section is deliberately outside the counted spans and is carried as standard prerequisites instead.\par\medskip
\noindent\textit{Editorial note.} The citation fields (source\_id, source\_version, authors, source\_url, license.name, source\_title) are identical to those of every other card this lane has produced from this chapter, so the batch selector treats them as one source.\par\medskip
\noindent\textit{Allowed standard prerequisites (editorial).}\begin{itemize}
\item Internal hom complexes and their compatibility with the monoidal structure. Remark \textbackslash{}ref\{remark-internal-hom\} and Categories, Remark \textbackslash{}ref\{categories-remark-internal-hom-monoidal\}, cited at proof lines 18176-18178. Produces the canonical map abstractly; the author then gives the direct construction to make the map explicit.
\item Sign rules for differentials on total complexes and internal hom complexes. used at proof lines 18181-18186, 18208-18213, 18252-18257, and 18268-18275. Determine the three-term normal forms of the differentials whose signs the construction has to match.
\end{itemize}\medskip
\ifdefined\hypertarget\hypertarget{sourcepointer-1}{}\fi
\noindent\textit{Source reference.} Original lemma 0A60: focused-seeds/AGP/sources/stacks-snapshot-20261011/tags.txt, source lines 12828--12828. The lemma's own canonical tag, quoted so that readers can resolve the label inside the delivered text; it points at the counted statement itself and carries no proof content.\par\medskip
\ifdefined\hypertarget\hypertarget{sourcepointer-2}{}\fi
\noindent\textit{Source reference.} Original lemma 0GWQ: focused-seeds/AGP/sources/stacks-snapshot-20261011/tags.txt, source lines 21080--21080. Navigation pointer to a statement the counted proof refers to but does not contain; nothing load-bearing is delegated to it.\par\medskip
\ifdefined\hypertarget\hypertarget{sourcepointer-3}{}\fi
\noindent\textit{Source reference.} Original lemma 0GWJ: focused-seeds/AGP/sources/stacks-snapshot-20261011/tags.txt, source lines 21074--21074. Navigation pointer to a statement the counted proof refers to but does not contain; nothing load-bearing is delegated to it.\par\medskip
\begin{lemma}
\label{lemma-evaluate-and-more}
Let $R$ be a ring. Given complexes $K^\bullet, L^\bullet, M^\bullet$
of $R$-modules there is a canonical morphism
$$
\text{Tot}(\Hom^\bullet(L^\bullet, M^\bullet) \otimes_R K^\bullet)
\longrightarrow
\Hom^\bullet(\Hom^\bullet(K^\bullet, L^\bullet), M^\bullet)
$$
of complexes of $R$-modules functorial in all three complexes.
\end{lemma}

\begin{proof}
Via the discussion in Remark \ref{remark-internal-hom} the existence
of such a canonical map follows from
Categories, Remark \ref{categories-remark-internal-hom-monoidal}.
We also give a direct construction.

\medskip\noindent
Consider an element $\beta$ of degree $n$ of the right hand side.
Then
$$
\beta = (\beta^{p, s}) \in
\prod\nolimits_{p + s = n} \Hom_R(\Hom^{-s}(K^\bullet, L^\bullet), M^p)
$$
Our sign rules tell us that
\begin{align*}
\text{d}(\beta^{p, s})
& =
\text{d}_M \circ \beta^{p, s}
- (-1)^n \beta^{p, s} \circ
\text{d}_{\Hom^\bullet(K^\bullet, L^\bullet)}
\end{align*}
We can describe the last term as follows
$$
(\beta^{p, s} \circ \text{d}_{\Hom^\bullet(K^\bullet, L^\bullet)})(f) =
\beta^{p, s}(\text{d}_L \circ f - (-1)^{s + 1} f \circ \text{d}_K)
$$
if $f \in \Hom^{-s - 1}(K^\bullet, L^\bullet)$. We conclude that in
some unspecified sense $\text{d}(\beta^{p, s})$ is a sum of three terms
with signs as follows
\begin{equation}
\label{equation-beta}
\text{d}(\beta^{p, s}) =
\text{d}_M(\beta^{p, s})
-(-1)^n\text{d}_L(\beta^{p, s}) +
(-1)^{p + 1}\text{d}_K(\beta^{p, s})
\end{equation}

\noindent
Next, we consider an element $\alpha$ of degree $n$ of the left hand side.
We can write it like so
$$
\alpha = (\alpha^{t, r}) \in
\bigoplus\nolimits_{t + r = n} \Hom^t(L^\bullet, M^\bullet) \otimes K^r
$$
Each $\alpha^{t, r}$ maps to an element
$$
\alpha^{t, r} \mapsto (\alpha^{p, q, r}) \in
\prod\nolimits_{p + q = t} \Hom_R(L^{-q}, M^p) \otimes_R K^r
$$
Our sign rules tell us that
\begin{align*}
\text{d}(\alpha^{p, q, r})
& =
\text{d}_{\Hom^\bullet(L^\bullet, M^\bullet)}(\alpha^{p, q, r})
+ (-1)^{p + q} \text{d}_K(\alpha^{p, q, r})
\end{align*}
where if we further write $\alpha^{p, q, r} = \sum g_i^{p, q} \otimes k_i^r$
then we have
$$
\text{d}_{\Hom^\bullet(L^\bullet, M^\bullet)}(\alpha^{p, q, r}) =
\sum (\text{d}_M \circ g_i^{p, q}) \otimes k_i^r
- (-1)^{p + q} \sum (g_i^{p, q} \circ \text{d}_L) \otimes k_i^r
$$
We conclude that in some unspecified sense $\text{d}(\alpha^{p, q, r})$
is a sum of three terms with signs as follows
\begin{equation}
\label{equation-alpha}
\text{d}(\alpha^{p, q, r}) =
\text{d}_M(\alpha^{p, q, r})
-(-1)^{p + q}\text{d}_L(\alpha^{p, q, r}) +
(-1)^{p + q}\text{d}_K(\alpha^{p, q, r})
\end{equation}

\noindent
To define our map we will use the canonical maps
$$
c_{p, q, r} : 
\Hom_R(L^{-q}, M^p) \otimes_R K^r
\longrightarrow
\Hom_R(\Hom_R(K^r, L^{-q}), M^p)
$$
which sends $\varphi \otimes k$ to the map $\psi \mapsto \varphi(\psi(k))$.
This is functorial in all three variables. With $s = q + r$
there is an inclusion
$$
\Hom_R(\Hom_R(K^r, L^{-q}), M^p) \subset
\Hom_R(\Hom^{-s}(K^\bullet, L^\bullet), M^p)
$$
coming from the projection
$\Hom^{-s}(K^\bullet, L^\bullet) \to \Hom_R(K^r, L^{-q})$.
Since $\alpha^{p, q, r}$
is nonzero only for a finite number of $r$ we see that for a given
$s$ there is only a finite number of $q, r$ with $q + r = s$.
Thus we can send $\alpha$ to the element $\beta$ with
$$
\beta^{p, s} =
\sum\nolimits_{q + r = s} \epsilon_{p, q, r} c_{p, q, r}(\alpha^{p, q, r})
$$
where where the sum uses the inclusions given above and where
$\epsilon_{p, q, r} \in \{\pm 1\}$. Comparing signs in the
equations (\ref{equation-beta}) and (\ref{equation-alpha})
we see that
\begin{enumerate}
\item $\epsilon_{p, q, r} = \epsilon_{p + 1, q, r}$
\item $-(-1)^n\epsilon_{p, q, r} = -(-1)^{p + q}\epsilon_{p, q - 1, r}$
or equivalently $\epsilon_{p, q, r} = (-1)^r\epsilon_{p, q - 1, r}$
\item $(-1)^{p + 1}\epsilon_{p, q, r} = (-1)^{p + q}\epsilon_{p, q, r + 1}$
or equivalently
$(-1)^{q + 1}\epsilon_{p, q, r} = \epsilon_{p, q, r + 1}$.
\end{enumerate}
A good solution is to take
$$
\epsilon_{p, r, s} = (-1)^{r + qr}
$$
The choice of this sign is explained in the remark following
the proof.
\end{proof}

\end{document}
